%% ============================================================
\section{Implementation and Technology Stack}
\label{sec:implementation}
%% ============================================================

\subsection{Project Structure}

\echoray{} is organised as a monorepo with two top-level directories:
\texttt{client/} and \texttt{server/}, each with independent
\texttt{package.json} manifests, \texttt{tsconfig.json} configurations, and
build pipelines.

\subsubsection{Server Structure}

The server follows an MVC pattern with a services layer for business logic:

\begin{verbatim}
server/src/
├── app.ts            — Express app setup (CORS, Helmet, Morgan)
├── server.ts         — HTTP server bootstrap + Socket.io init
├── socket.ts         — Real-time event handlers
├── controllers/      — HTTP request handlers (thin layer)
│   ├── user.controller.ts
│   ├── project.controller.ts
│   └── ai.controller.ts
├── services/         — Business logic
│   ├── user.service.ts
│   ├── project.service.ts
│   ├── ai.service.ts
│   └── redis.service.ts
├── routes/           — Express router definitions
├── middlewares/      — auth.middleware.ts (JWT + Redis check)
├── db/index.ts       — Prisma client (singleton + Accelerate)
└── utils/            — auth.ts, schema.ts, type.ts
\end{verbatim}

\subsubsection{Client Structure}

The client uses the Next.js App Router introduced in Next.js~13:

\begin{verbatim}
client/
├── app/
│   ├── page.tsx                — Root redirect to /home
│   ├── layout.tsx              — Root layout (ReduxProvider)
│   ├── signin/page.tsx         — Sign-in form
│   ├── signup/page.tsx         — Registration form
│   ├── logout/page.tsx         — Logout + redirect
│   ├── home/page.tsx           — Project dashboard
│   └── project/[id]/page.tsx   — Collaborative workspace (718 LOC)
├── components/
│   ├── CodeEditor.tsx          — Monaco Editor wrapper
│   ├── ChatBox.tsx             — Real-time chat UI
│   ├── ProjectFiles.tsx        — Virtual file tree navigator
│   ├── IframePreview.tsx       — WebContainer live preview
│   ├── CollaboratorList.tsx    — Active member display
│   ├── AddCollaboratorModal.tsx— Collaborator management
│   └── UserProtectWrapper.tsx  — Auth guard HOC
├── config/
│   ├── socketIo.ts             — Socket.io client singleton
│   └── Webcontainer.ts         — WebContainer boot singleton
└── store/
    ├── index.ts                — Redux store configuration
    ├── userSlice.ts            — User state reducer
    └── ReduxProvider.tsx       — Provider wrapper
\end{verbatim}

\subsection{Key Technology Decisions}

\subsubsection{Next.js 15 with App Router}

Next.js~15 provides hybrid SSR/SSG rendering, file-system routing, and
React~19 Server Components. For \echoray{}, the primary benefit is the
App Router's support for nested layouts, enabling the project workspace
(\texttt{project/[id]/page.tsx}) to be a rich client-side component while
authentication pages benefit from server-side rendering for SEO and
first-paint performance.

\subsubsection{Monaco Editor}

Monaco~\cite{monacoeditor2023}, the editor powering Visual Studio Code, is
embedded via the \texttt{@monaco-editor/react} package. It provides:

\begin{itemize}
  \item Syntax highlighting for 60+ languages via TextMate grammars.
  \item IntelliSense code completion backed by the TypeScript Language Server.
  \item Multi-cursor editing, minimap, and diff views.
  \item A programmatic API for reading/writing content, enabling integration
    with \echoray{}'s file-tree state.
\end{itemize}

\subsubsection{Prisma ORM with Accelerate}

Prisma generates fully typed query methods from the schema, eliminating
runtime type errors at the database boundary. Prisma Accelerate adds:

\begin{itemize}
  \item A globally distributed connection pool (reducing cold-start latency
    from serverless functions).
  \item An LRU query cache with configurable TTLs, reducing median read
    latency by~$\sim$73\% in our measurements (\cref{sec:evaluation}).
\end{itemize}

\subsubsection{TailwindCSS 4 and GSAP}

TailwindCSS~4 introduces a JIT compiler and CSS variables-based theme system,
reducing the need for custom CSS. GSAP (GreenSock Animation Platform) provides
60\,fps declarative animations for page transitions and micro-interactions,
improving perceived responsiveness.

\subsubsection{Redux Toolkit}

Redux Toolkit~\cite{reduxToolkit2023} provides a \emph{slice}-based API that
collapses the boilerplate of traditional Redux (separate action creators,
reducers, and thunks). The \texttt{userSlice} (\texttt{store/userSlice.ts})
manages a single piece of state:

\begin{lstlisting}[caption={Redux user slice.}, label=lst:redux]
import { createSlice, PayloadAction } from "@reduxjs/toolkit";

interface UserState { user: User | null; }

const userSlice = createSlice({
  name: "user",
  initialState: { user: null } as UserState,
  reducers: {
    setUser: (state, action: PayloadAction<User>) => {
      state.user = action.payload;
    },
    clearUser: (state) => {
      state.user = null;
    },
  },
});

export const { setUser, clearUser } = userSlice.actions;
export default userSlice.reducer;
\end{lstlisting}

\subsection{Build and Deployment Pipeline}

\Cref{tab:buildpipeline} summarises the build pipeline steps for both
front end and back end.

\begin{table}[H]
  \centering
  \caption{Build and deployment pipeline summary.}
  \label{tab:buildpipeline}
  \begin{tabular}{lll}
    \toprule
    \textbf{Stage}     & \textbf{Frontend}           & \textbf{Backend} \\
    \midrule
    Type check         & \texttt{tsc --noEmit}       & \texttt{tsc --noEmit} \\
    Build              & \texttt{next build}         & \texttt{tsc -p tsconfig.json} \\
    Output             & \texttt{.next/} (SSR bundle) & \texttt{dist/} (JS) \\
    Dev server         & \texttt{next dev} (port 3000)& \texttt{ts-node-dev} (port 5000) \\
    Deploy target      & Vercel (auto-deploy on push) & Render/Railway \\
    DB migration       & N/A                         & \texttt{prisma migrate deploy} \\
    \bottomrule
  \end{tabular}
\end{table}

\subsection{Configuration and Environment Variables}

\Cref{tab:envvars} lists all required environment variables.

\begin{table}[H]
  \centering
  \caption{Required environment variables for \echoray{} deployment.}
  \label{tab:envvars}
  \begin{tabular}{lll}
    \toprule
    \textbf{Variable}         & \textbf{Location} & \textbf{Purpose} \\
    \midrule
    \texttt{DATABASE\_URL}    & server            & PostgreSQL connection string \\
    \texttt{JWT\_SECRET}      & server            & HMAC-SHA256 signing key \\
    \texttt{AI\_API\_KEY}     & server            & Google Generative AI API key \\
    \texttt{REDIS\_URL}       & server            & Redis connection URL \\
    \texttt{NEXT\_PUBLIC\_API\_URL} & client     & Backend base URL \\
    \bottomrule
  \end{tabular}
\end{table}

\subsection{Code Metrics}

\Cref{tab:codemetrics} provides a quantitative overview of the codebase.

\begin{table}[H]
  \centering
  \caption{Codebase metrics for \echoray{}.}
  \label{tab:codemetrics}
  \begin{tabular}{lrr}
    \toprule
    \textbf{Metric}            & \textbf{Frontend} & \textbf{Backend} \\
    \midrule
    Source files (.ts/.tsx)    & 25                & 23 \\
    Lines of code (approx.)    & 1\,500+           & 800+ \\
    Database models            & ---               & 2 \\
    REST endpoints             & ---               & 14 \\
    Socket.io event types      & 1 (bidirectional) & 1 \\
    Redux slices               & 1                 & --- \\
    Zod validation schemas     & ---               & 6 \\
    \bottomrule
  \end{tabular}
\end{table}
